\section{Conclusion}

The breathable Moon combines conditions that make an Earth model's inherited values consequential in several
parts of the same calculation. At a familiar surface pressure, weak gravity increases the mass of absorbing
gas and stretches the air through a substantial fraction of the radius, while the long solar day changes the
duration of heating. Adapting the climate, storm and wave codes to this combination identified both physical
approximations requiring new evidence and defects whose effects extend, by the code's arithmetic, to
configurations used for natural planets. A design developed for one world thereby supplied tests relevant to
the modelling of others.

The three ExoPlaSim defects illustrate that connection particularly clearly. Their errors grew with the test
world's long day and small tilt, although the code continued to compute climates. Two redistributed sunlight
while preserving its global mean at the top of the air; correcting them warmed the surface by less than a
kelvin over a two-year test, still gaining heat, while the fraction of land receiving little rain grew
fourfold. The altered sunlight changed the clouds and their reflection, so its distribution affected both
the energy balance and where water reached the ground. For the slowly rotating planets on which these
defects also act, the corresponding climate change remains to be measured. The lunar result makes that a
question about the spatial organisation of the climate as well as its mean warmth.

A similar distinction between an integrated quantity and its distribution arose within the radiation.
PlaSim could be calibrated to an outgoing flux while cooling the wrong layers; its water absorption could
be fitted to the column's total while taking up sunlight that carbon dioxide and oxygen should have absorbed
higher in the air. Those differences matter to the comparison with resolved convection: the GCM determines
the heating and circulation from which CM1's storms develop, and CM1 computes its own radiative cooling
within that imposed environment. A disagreement between the models can therefore originate in the vertical
distribution of energy before their convection schemes respond to it.

At the ground, that distribution becomes a question about the depth through which heat is mixed. In the
lunar comparison, imposing the transport across the ring removed part of CM1's warm offset, while comparison
with the radiation reference identified excess low-level cooling in both models. The remainder includes
their different treatment of upward heat transport, with the starting sounding's high ice also contributing
an unquantified term. The two calculations consequently supply a set of physical explanations to
discriminate. A warmer resolved atmosphere becomes evidence about the GCM's convection only after the
forcing and radiative differences have been accounted for; the choice of cloud calibration depends on the
same separation.

The proposed use of resolved cloud to calibrate the GCM thus entails agreement on the environment in which
that cloud develops. With consistent radiative heating and the missing large-scale transport imposed,
CM1's condensate and cover could constrain the GCM's cloud relations at the world's own gravity and day.
That calibration would feed back into the global heating and circulation, requiring a new comparison with
the resolved runs. The resulting iteration would test whether the two descriptions of cloud can agree in
one climate, with the resolved model's microphysics determining how much confidence to place in the target.

Rain makes that last dependence consequential for the surface. In the lunar box, most of the water leaving
the clouds evaporated during its descent, and the ground received about a quarter of the forcing GCM's rain.
The fall-speed repair follows gravity through weight and drag, while the limits on drop size and breakup
remain fitted to Earth. Larger drops could survive more of the descent; their possible size increases with
the capillary length, although collisions determine which sizes actually occur. Testing the retained size
limits would measure how much of the rain gap comes from that choice within CM1. Only after this
discrimination could the smaller resolved rainfall be used to revise the water budgets of the proposed lakes.

The need for that test follows from the different scalings within a storm. The dry dynamics stretch in
height and time as gravity falls, while the fall speeds of droplets, rain and graupel change by different
powers. A stretched terrestrial storm therefore has different ratios of ascent, sedimentation and
evaporation times. The wave calculations offer a useful comparison: on the paths tested, their gravity
scaling preserves the sea's dimensionless form with Earth's growth and breaking coefficients retained.
The storm's particle processes depart from the dry similarity, so their effect on cloud and precipitation
has to be computed. Changing gravity can alter how water passes through a storm as well as the height at
which the storm develops.

Some of the particle physics can be tested in the conditions of an Earth laboratory. In the charging
calculations, two laws fitted to laboratory measurements produced different arrangements of charge in
windows of the lunar runs because they assigned different signs to rimed ice in little cloud water.
Measurements at the simulated collision speeds, temperatures and water contents could distinguish those
laws for the relevant particles. They would constrain the electrical structure of a storm tens of
kilometres high through collisions reproduced at laboratory scale. Ground-strike rates would require the
additional calculation of a leader's descent through the deep air. The observable process and the
planetary consequence can be investigated at different scales.

That separation gives a practical route to knowledge of a world whose climate is unobserved. The gas
absorption can be calculated from spectroscopy, the charging laws compared with measurements at the
collision conditions, and the scaled wave dynamics checked against similarity. The remaining dependence
on empirical coefficients stays specific to the process: cloud microphysics in the rain, charge transfer
in the electrical structure, growth and breaking in the sea. A calculation of a new world can then
identify the measurement or comparison on which a particular feature depends. The lake's rain supply and
the storm's charge arrangement become distinct research questions with distinct possible outcomes.

For a constructed world, a result of such a comparison can also change a chosen boundary condition.
Correcting the lunar radiation changed the shield's dimming required to maintain the design temperature,
and a consistent longwave fit or a cloud calibration could change it again. The uncertain quantity is the
predicted heating; the available design control is the illumination. Holding the target temperature while
revising that control would require the circulation, cloud and rain to be recomputed, with new forcing for
the storm and sea models. Two versions of the design could satisfy the same temperature target while
requiring different shielding and supplying different amounts of water to the land.

For a natural planet, the corresponding calculation would hold the prescribed starlight and determine how
far the revised physics moves the climate. This gives the two applications a shared experimental basis:
a process constrained on the designed world can be tested under a natural planet's conditions, while a
comparison between models on that planet can test a proposed use of the repaired code. The general repairs
are built to transfer within their laws; the geometry needs the new world's column, and empirical
calibrations need evidence over its conditions. Beyond Earth's values and numerical controls, the
adapted models have so far run on the lunar case alone. Further worlds would test the interactions of
processes that the lunar adaptation has made individually accessible.

For the radiation, a checked scheme reduced from the reference spectroscopy would make such comparisons
possible across its gases, pressures and temperatures under the starlight for which it was built.
For cloud, comparisons with resolved convection would supply a calibration for the world being studied.
Together they would let a modeller follow a change in illumination through absorption, convection and
cloud into the water arriving at the ground, while the remaining particle laws were tested separately.
The scientific prospect is to determine how the same air and water organise into different environments
as planetary conditions change, including differences that remain small in a global temperature.

Those environments extend beyond the columns in which their radiation is calculated. The wind that moves
heat and moisture also forces a sea, whose spectra determine the waves arriving at a particular shore;
rainfall in the surrounding terrain determines whether a basin can retain the lake assumed in the climate.
Following these dependencies would bring the atmosphere, water distribution and coastal conditions of a
proposed world into the same physical calculation. A designer could then compare changes to illumination
or the distribution of surface water through their effects on a place, while retaining the body's own
gravity, size and day.

On a small world with breathable air, those native conditions can produce tall storms above slowly falling
rain and long-period waves under light winds. Its upper air can remain sunlit after the ground has entered
a night that lasts many Earth days. As the relations between these processes are established, a world can
be considered through the environments it could sustain: a basin replenished through the long day, a shore
under the swell generated by its winds, and the changing light and weather through which life would grow.
The purpose of extending these models is to make such environments open to physical investigation, so that
the exploration of possible worlds can inform both the interpretation of planets around other stars and
the choice of conditions for worlds people may one day inhabit.
